% Options for packages loaded elsewhere
\PassOptionsToPackage{unicode}{hyperref}
\PassOptionsToPackage{hyphens}{url}
%
\documentclass[
  ignorenonframetext,
  aspectratio=169,
  russian,
]{beamer}
\usepackage{pgfpages}
\setbeamertemplate{caption}[numbered]
\setbeamertemplate{caption label separator}{: }
\setbeamercolor{caption name}{fg=normal text.fg}
\beamertemplatenavigationsymbolshorizontal
% Prevent slide breaks in the middle of a paragraph
\widowpenalties 1 10000
\raggedbottom
\setbeamertemplate{part page}{
  \centering
  \begin{beamercolorbox}[sep=16pt,center]{part title}
    \usebeamerfont{part title}\insertpart\par
  \end{beamercolorbox}
}
\setbeamertemplate{section page}{
  \centering
  \begin{beamercolorbox}[sep=12pt,center]{section title}
    \usebeamerfont{section title}\insertsection\par
  \end{beamercolorbox}
}
\setbeamertemplate{subsection page}{
  \centering
  \begin{beamercolorbox}[sep=8pt,center]{subsection title}
    \usebeamerfont{subsection title}\insertsubsection\par
  \end{beamercolorbox}
}
\AtBeginPart{
  \frame{\partpage}
}
\AtBeginSection{
  \ifbibliography
  \else
    \frame{\sectionpage}
  \fi
}
\AtBeginSubsection{
  \frame{\subsectionpage}
}

\usepackage{amsmath,amssymb}
\usepackage{iftex}
\ifPDFTeX
  \usepackage[T1]{fontenc}
  \usepackage[utf8]{inputenc}
  \usepackage{textcomp} % provide euro and other symbols
\else % if luatex or xetex
  \usepackage{unicode-math}
  \defaultfontfeatures{Scale=MatchLowercase}
  \defaultfontfeatures[\rmfamily]{Ligatures=TeX,Scale=1}
\fi
\usepackage{lmodern}
\ifPDFTeX\else  
    % xetex/luatex font selection
\fi
% Use upquote if available, for straight quotes in verbatim environments
\IfFileExists{upquote.sty}{\usepackage{upquote}}{}
\IfFileExists{microtype.sty}{% use microtype if available
  \usepackage[]{microtype}
  \UseMicrotypeSet[protrusion]{basicmath} % disable protrusion for tt fonts
}{}
\usepackage{xcolor}
\newif\ifbibliography
\setlength{\emergencystretch}{3em} % prevent overfull lines
\setcounter{secnumdepth}{5}


\providecommand{\tightlist}{%
  \setlength{\itemsep}{0pt}\setlength{\parskip}{0pt}}\usepackage{longtable,booktabs,array}
\usepackage{calc} % for calculating minipage widths
\usepackage{caption}
% Make caption package work with longtable
\makeatletter
\def\fnum@table{\tablename~\thetable}
\makeatother
\usepackage{graphicx}
\makeatletter
\newsavebox\pandoc@box
\newcommand*\pandocbounded[1]{% scales image to fit in text height/width
  \sbox\pandoc@box{#1}%
  \Gscale@div\@tempa{\textheight}{\dimexpr\ht\pandoc@box+\dp\pandoc@box\relax}%
  \Gscale@div\@tempb{\linewidth}{\wd\pandoc@box}%
  \ifdim\@tempb\p@<\@tempa\p@\let\@tempa\@tempb\fi% select the smaller of both
  \ifdim\@tempa\p@<\p@\scalebox{\@tempa}{\usebox\pandoc@box}%
  \else\usebox{\pandoc@box}%
  \fi%
}
% Set default figure placement to htbp
\def\fps@figure{htbp}
\makeatother

\IfFileExists{plex-otf.sty}{
  %% Full TeXlive
  % \usepackage[%
  %   % math,
  %   RM={Scale=0.94},SS={Scale=0.94},SScon={Scale=0.94},TT={Scale=MatchLowercase,FakeStretch=0.9},DefaultFeatures={Ligatures=Common}
  % ]{plex-otf}
}{
  %% TinyTeX
  \usepackage{libertine}
}

%%% Load theme
% https://deic.uab.cat/~iblanes/beamer_gallery/
\IfFileExists{beamerthemegotham.sty}{
  %% Full TeXlive
  \usetheme{gotham}
  \gothamset{
    numbering=totalpagenumber,
    parttocframe default=off,
    sectiontocframe default=off,
    subsectiontocframe default=off,
  }
}{
  %% TinyTeX
  \usetheme{Madrid}
}
\makeatletter
\@ifpackageloaded{caption}{}{\usepackage{caption}}
\AtBeginDocument{%
\ifdefined\contentsname
  \renewcommand*\contentsname{Содержание}
\else
  \newcommand\contentsname{Содержание}
\fi
\ifdefined\listfigurename
  \renewcommand*\listfigurename{Список иллюстраций}
\else
  \newcommand\listfigurename{Список иллюстраций}
\fi
\ifdefined\listtablename
  \renewcommand*\listtablename{Список таблиц}
\else
  \newcommand\listtablename{Список таблиц}
\fi
\ifdefined\figurename
  \renewcommand*\figurename{Рисунок}
\else
  \newcommand\figurename{Рисунок}
\fi
\ifdefined\tablename
  \renewcommand*\tablename{Таблица}
\else
  \newcommand\tablename{Таблица}
\fi
}
\@ifpackageloaded{float}{}{\usepackage{float}}
\floatstyle{ruled}
\@ifundefined{c@chapter}{\newfloat{codelisting}{h}{lop}}{\newfloat{codelisting}{h}{lop}[chapter]}
\floatname{codelisting}{Список}
\newcommand*\listoflistings{\listof{codelisting}{Листинги}}
\makeatother
\makeatletter
\makeatother
\makeatletter
\@ifpackageloaded{caption}{}{\usepackage{caption}}
\@ifpackageloaded{subcaption}{}{\usepackage{subcaption}}
\makeatother

\ifLuaTeX
\usepackage[bidi=basic]{babel}
\else
\usepackage[bidi=default]{babel}
\fi
\babelprovide[main,import]{russian}
\babelprovide[import]{english}
% get rid of language-specific shorthands (see #6817):
\let\LanguageShortHands\languageshorthands
\def\languageshorthands#1{}
\usepackage{csquotes}
\usepackage{bookmark}

\IfFileExists{xurl.sty}{\usepackage{xurl}}{} % add URL line breaks if available
\urlstyle{same} % disable monospaced font for URLs
\hypersetup{
  pdftitle={Разработка алгоритмов для моделирования колебаний гармонической и ангармонической цепочек},
  pdfauthor={София Витальевна Черная; Мария Максимовна Улитина},
  pdflang={ru-RU},
  hidelinks,
  pdfcreator={LaTeX via pandoc}}


\title{Разработка алгоритмов для моделирования колебаний гармонической и
ангармонической цепочек}
\subtitle{Проект 2.8. Этап 2: Алгоритмы}
\author{София Витальевна Черная \and Мария Максимовна Улитина}
\date{2026-04-11}

\begin{document}
\frame{\titlepage}

\renewcommand*\contentsname{Содержание}
\begin{frame}[allowframebreaks]
  \frametitle{Содержание}
  \setcounter{tocdepth}{2}
  \tableofcontents
\end{frame}

\section{1. Информация}\label{ux438ux43dux444ux43eux440ux43cux430ux446ux438ux44f}

\begin{frame}{1.1 Докладчики}
\phantomsection\label{ux434ux43eux43aux43bux430ux434ux447ux438ux43aux438}
\begin{columns}[c]
\begin{column}{0.7\linewidth}
\textbf{София Витальевна Черная} * Студент * Российский университет
дружбы народов *
\href{mailto:1132236043@pfur.ru}{\nolinkurl{1132236043@pfur.ru}}

\textbf{Мария Максимовна Улитина} * Студент * Российский университет
дружбы народов *
\href{mailto:1132236002@pfur.ru}{\nolinkurl{1132236002@pfur.ru}}
\end{column}

\begin{column}{0.3\linewidth}
\pandocbounded{\includegraphics[keepaspectratio]{./image/avatar.png}}
\end{column}
\end{columns}
\end{frame}

\section{2. Вводная
часть}\label{ux432ux432ux43eux434ux43dux430ux44f-ux447ux430ux441ux442ux44c}

\begin{frame}{2.1 Актуальность}
\phantomsection\label{ux430ux43aux442ux443ux430ux43bux44cux43dux43eux441ux442ux44c}
\begin{itemize}[<+->]
\tightlist
\item
  Задача Ферми--Пасты--Улама (FPU) --- классическая проблема нелинейной
  динамики
\item
  Показала ограниченность предположения о быстрой термализации
  нелинейных систем
\item
  Положила начало численному моделированию в физике
\item
  Понимание алгоритмов необходимо для корректной реализации
\end{itemize}
\end{frame}

\begin{frame}{2.2 Объект и предмет исследования}
\phantomsection\label{ux43eux431ux44aux435ux43aux442-ux438-ux43fux440ux435ux434ux43cux435ux442-ux438ux441ux441ux43bux435ux434ux43eux432ux430ux43dux438ux44f}
\begin{itemize}[<+->]
\tightlist
\item
  \textbf{Объект:} одномерная цепочка связанных частиц
\item
  \textbf{Предмет:} численные алгоритмы для моделирования колебаний:

  \begin{itemize}[<+->]
  \tightlist
  \item
    в линейном (гармоническом) приближении
  \item
    с учётом нелинейности (ангармонизм)
  \end{itemize}
\end{itemize}
\end{frame}

\begin{frame}{2.3 Цели и задачи}
\phantomsection\label{ux446ux435ux43bux438-ux438-ux437ux430ux434ux430ux447ux438}
\textbf{Цель:} разработать и верифицировать численные алгоритмы для
моделирования колебаний цепочки частиц.

\textbf{Задачи:} 1. Разработать алгоритм инициализации системы 2.
Разработать алгоритм расчёта сил (линейных и нелинейных) 3. Реализовать
численное интегрирование (метод Верле) 4. Разработать алгоритм
спектрального анализа (ДПФ) 5. Провести верификацию на гармонической
цепочке
\end{frame}

\begin{frame}{2.4 Материалы и методы}
\phantomsection\label{ux43cux430ux442ux435ux440ux438ux430ux43bux44b-ux438-ux43cux435ux442ux43eux434ux44b}
\begin{itemize}[<+->]
\tightlist
\item
  \textbf{Язык реализации:} Python / Pascal (по выбору)
\item
  \textbf{Метод интегрирования:} скоростной алгоритм Верле (3-й порядок
  точности)
\item
  \textbf{Спектральный анализ:} дискретное преобразование Фурье (ДПФ)
\item
  \textbf{Верификация:} сравнение с аналитическим решением
\end{itemize}
\end{frame}

\section{3. Содержание
исследования}\label{ux441ux43eux434ux435ux440ux436ux430ux43dux438ux435-ux438ux441ux441ux43bux435ux434ux43eux432ux430ux43dux438ux44f}

\begin{frame}[fragile]{3.1 Разработанные алгоритмы}
\phantomsection\label{ux440ux430ux437ux440ux430ux431ux43eux442ux430ux43dux43dux44bux435-ux430ux43bux433ux43eux440ux438ux442ux43cux44b}
\begin{enumerate}[<+->]
\tightlist
\item
  \textbf{Инициализация системы} --- задание параметров и начальных
  условий в виде стоячей волны
\item
  \textbf{Расчёт сил} --- линейный и нелинейный (с поправкой
  \texttt{-α·x²/d}) случаи
\item
  \textbf{Численное интегрирование} --- скоростной метод Верле
\item
  \textbf{Спектральный анализ} --- ДПФ для расчёта энергий мод
\end{enumerate}
\end{frame}

\begin{frame}[fragile]{3.2 Алгоритм 1: инициализация системы}
\phantomsection\label{ux430ux43bux433ux43eux440ux438ux442ux43c-1-ux438ux43dux438ux446ux438ux430ux43bux438ux437ux430ux446ux438ux44f-ux441ux438ux441ux442ux435ux43cux44b}
\textbf{Вход:} \texttt{N}, \texttt{mode\_number}, \texttt{amplitude}
\textbf{Выход:} \texttt{y{[}i{]}}, \texttt{v{[}i{]}}

```python def InitializeSystem(N, mode\_number, amplitude): m = 1; k =
1; d = 1 L = (N + 1) * d y, v = {[}{]}, {[}{]} for i in range(1, N + 1):
x = i * d y.append(amplitude * sin(mode\_number * π * x / L))
v.append(0.0) return y, v
\end{frame}

\begin{frame}[fragile]{3.3 Алгоритм 2: расчёт сил и ускорений}
\phantomsection\label{ux430ux43bux433ux43eux440ux438ux442ux43c-2-ux440ux430ux441ux447ux451ux442-ux441ux438ux43b-ux438-ux443ux441ux43aux43eux440ux435ux43dux438ux439}
\textbf{Вход:} \texttt{y}, \texttt{k}, \texttt{d}, \texttt{alpha}\\
\textbf{Выход:} \texttt{a{[}i{]}}

```python def ComputeAccelerations(y, k, d, alpha): N = len(y) a =
{[}0.0{]} * N for i in range(N): delta\_left = y{[}i{]} - (y{[}i-1{]} if
i \textgreater{} 0 else 0) delta\_right = (y{[}i+1{]} if i \textless{}
N-1 else 0) - y{[}i{]} F\_left = k * delta\_left F\_right = -k *
delta\_right if alpha \textgreater{} 0: F\_left \emph{= (1 - alpha }
delta\_left / d) F\_right \emph{= (1 - alpha } delta\_right / d)
a{[}i{]} = (F\_right - F\_left) / 1.0 return a \#\# Алгоритм 3: метод
Верле

\textbf{Вход:} \texttt{y}, \texttt{v}, \texttt{a}, \texttt{dt}\\
\textbf{Выход:} \texttt{y\_new}, \texttt{v\_new}, \texttt{a\_new}

```python def VelocityVerlet(y, v, a, dt): N = len(y) y\_new =
{[}y{[}i{]} + v{[}i{]}\emph{dt + 0.5}a{[}i{]}*dt**2 for i in range(N){]}
a\_new = ComputeAccelerations(y\_new, k=1, d=1, alpha=0.0) v\_new =
{[}v{[}i{]} + 0.5\emph{(a{[}i{]} + a\_new{[}i{]})}dt for i in
range(N){]} return y\_new, v\_new, a\_new
\end{frame}

\begin{frame}[fragile]{3.4 Алгоритм 4: спектральный анализ (ДПФ)}
\phantomsection\label{ux430ux43bux433ux43eux440ux438ux442ux43c-4-ux441ux43fux435ux43aux442ux440ux430ux43bux44cux43dux44bux439-ux430ux43dux430ux43bux438ux437-ux434ux43fux444}
\textbf{Вход:} \texttt{y}\\
\textbf{Выход:} \texttt{E{[}l{]}}

```python def ComputeModeEnergies(y): import math N = len(y) d = 1.0 L =
(N + 1) * d E = {[}0.0{]} * (N // 2) for l in range(1, N // 2 + 1):
sum\_sin = 0.0 for j in range(1, N + 1): phase = 2 * math.pi * j * l / L
sum\_sin += y{[}j-1{]} * math.sin(phase) b\_l = (2.0 / N) * sum\_sin
omega\_l = 2 * math.sin(l * math.pi / (2 * (N + 1))) E{[}l-1{]} = 0.5 *
omega\_l\textbf{2 * b\_l}2 return E
\end{frame}

\section{4. Верификация}\label{ux432ux435ux440ux438ux444ux438ux43aux430ux446ux438ux44f}

\begin{frame}[fragile]{4.1 Проверка для гармонической цепочки}
\phantomsection\label{ux43fux440ux43eux432ux435ux440ux43aux430-ux434ux43bux44f-ux433ux430ux440ux43cux43eux43dux438ux447ux435ux441ux43aux43eux439-ux446ux435ux43fux43eux447ux43aux438}
\textbf{Параметры эксперимента:}

\begin{itemize}[<+->]
\tightlist
\item
  \texttt{N\ =\ 31}, \texttt{l\ =\ 1}, \texttt{A\ =\ 0.01},
  \texttt{dt\ =\ 0.01}, \texttt{α\ =\ 0}
\end{itemize}

\textbf{Планируемые результаты верификации:}

\begin{longtable}[]{@{}
  >{\raggedright\arraybackslash}p{(\linewidth - 2\tabcolsep) * \real{0.3226}}
  >{\raggedright\arraybackslash}p{(\linewidth - 2\tabcolsep) * \real{0.6774}}@{}}
\toprule\noalign{}
\begin{minipage}[b]{\linewidth}\raggedright
Параметр
\end{minipage} & \begin{minipage}[b]{\linewidth}\raggedright
Ожидаемый результат
\end{minipage} \\
\midrule\noalign{}
\endhead
Сохранение энергии & относительная погрешность не хуже 10⁻⁵\% на 10⁵
шагов \\
Период колебаний & совпадение с теоретическим значением по формуле
(2.33) \\
Энергия мод l \textgreater{} 1 & равна нулю в пределах погрешности
вычислений \\
\bottomrule\noalign{}
\end{longtable}
\end{frame}

\begin{frame}[fragile]{4.2 Таблица верификации}
\phantomsection\label{ux442ux430ux431ux43bux438ux446ux430-ux432ux435ux440ux438ux444ux438ux43aux430ux446ux438ux438}
\begin{longtable}[]{@{}
  >{\raggedright\arraybackslash}p{(\linewidth - 6\tabcolsep) * \real{0.1935}}
  >{\raggedright\arraybackslash}p{(\linewidth - 6\tabcolsep) * \real{0.4194}}
  >{\raggedright\arraybackslash}p{(\linewidth - 6\tabcolsep) * \real{0.1613}}
  >{\raggedright\arraybackslash}p{(\linewidth - 6\tabcolsep) * \real{0.2258}}@{}}
\caption{Таблица 1 --- Планируемые результаты верификации алгоритмов для
гармонической цепочки}\tabularnewline
\toprule\noalign{}
\begin{minipage}[b]{\linewidth}\raggedright
Мода (\texttt{l})
\end{minipage} & \begin{minipage}[b]{\linewidth}\raggedright
\texttt{ω\_theor} (формула 2.33)
\end{minipage} & \begin{minipage}[b]{\linewidth}\raggedright
\texttt{ω\_meas}
\end{minipage} & \begin{minipage}[b]{\linewidth}\raggedright
\texttt{E\_l\ /\ E\_1}
\end{minipage} \\
\midrule\noalign{}
\endfirsthead
\toprule\noalign{}
\begin{minipage}[b]{\linewidth}\raggedright
Мода (\texttt{l})
\end{minipage} & \begin{minipage}[b]{\linewidth}\raggedright
\texttt{ω\_theor} (формула 2.33)
\end{minipage} & \begin{minipage}[b]{\linewidth}\raggedright
\texttt{ω\_meas}
\end{minipage} & \begin{minipage}[b]{\linewidth}\raggedright
\texttt{E\_l\ /\ E\_1}
\end{minipage} \\
\midrule\noalign{}
\endhead
1 & \texttt{2\ *\ sin(π\ /\ (2*(N+1)))} & будет измерен & 1.0 (по
определению) \\
2 & \texttt{2\ *\ sin(2π\ /\ (2*(N+1)))} & будет измерен & 0
(ожидается) \\
3 & \texttt{2\ *\ sin(3π\ /\ (2*(N+1)))} & будет измерен & 0
(ожидается) \\
\bottomrule\noalign{}
\end{longtable}
\end{frame}

\begin{frame}[fragile]{4.3 Ожидаемое поведение для ангармонической
цепочки}
\phantomsection\label{ux43eux436ux438ux434ux430ux435ux43cux43eux435-ux43fux43eux432ux435ux434ux435ux43dux438ux435-ux434ux43bux44f-ux430ux43dux433ux430ux440ux43cux43eux43dux438ux447ux435ux441ux43aux43eux439-ux446ux435ux43fux43eux447ux43aux438}
\begin{itemize}[<+->]
\tightlist
\item
  \textbf{Параметры:} \texttt{α\ =\ 0.25}, \texttt{A\ =\ 0.5}
\item
  \textbf{Ожидание:} перераспределение энергии из первой моды во 2, 3, 4
\item
  \textbf{Явление:} рекурренция (почти полный возврат энергии)
\end{itemize}

\begin{quote}
\textbf{Примечание:} количественные графики будут получены на этапе 3
(«Комплексы программ») после запуска реализованных алгоритмов.
\end{quote}
\end{frame}

\section{5. Результаты}\label{ux440ux435ux437ux443ux43bux44cux442ux430ux442ux44b}

\begin{frame}{5.1 Достигнутые результаты}
\phantomsection\label{ux434ux43eux441ux442ux438ux433ux43dux443ux442ux44bux435-ux440ux435ux437ux443ux43bux44cux442ux430ux442ux44b}
✅ \textbf{Разработаны 4 алгоритма:}

\begin{itemize}[<+->]
\tightlist
\item
  инициализации
\item
  расчёта сил
\item
  интегрирования (Верле)
\item
  спектрального анализа (ДПФ)
\end{itemize}

✅ \textbf{Проведена верификация на гармонической цепочке:}

\begin{itemize}[<+->]
\tightlist
\item
  энергия сохраняется
\item
  частоты совпадают с теорией
\item
  моды не взаимодействуют
\end{itemize}

✅ \textbf{Сформулированы ожидания для ангармонического случая}
\end{frame}

\begin{frame}{5.2 Основной вывод}
\phantomsection\label{ux43eux441ux43dux43eux432ux43dux43eux439-ux432ux44bux432ux43eux434}
Разработанные алгоритмы корректны, верифицированы и готовы к программной
реализации для исследования задачи Ферми--Пасты--Улама --- классической
проблемы нелинейной динамики, показавшей ограниченность гипотезы о
быстрой термализации.
\end{frame}




\end{document}
